% ============================================================================
\clearpage
\section*{v4.9.0-r79 Source-Only: Gauge Independence of the Four Hidden Conditions of the OP5 Comparison System}
\addcontentsline{toc}{section}{v4.9.0-r79 Source-Only: Gauge Independence of the Four Hidden Conditions of the OP5 Comparison System}
\label{sec:v490-r79-gauge-independence}
% ============================================================================

\paragraph{Purpose.}
Layer v4.9.0-r78 showed that \(R\) and \(\Omega\) are covariant, while the gauge \(\Psi^\dagger K_2\Psi=S\) keeps only
\(\mathfrak{so}(1,2)\) (\S\ref{sec:v490-r78-value-covariance}).  The four hidden conditions come from \(R\) together with this
gauge, so their status was open: they could be content of the covariant system or an artefact of the non-covariant gauge
fixing.  On 4 November 2026 the author asked to decide this.  This layer records one pre-registered test (H67), which is
confirmed.  Everything concerns the comparison model of \S\ref{sec:v410-state-mediator} only.  Earlier layers are not
rewritten.

\noindent\textbf{Source-only status.}
This v4.9.0-r79 layer is source-only.  It does not build a PDF, does not assign a new DOI,
and does not authorize a public release.  The v4.9.0 Version DOI
\href{https://doi.org/10.5281/zenodo.23534954}{10.5281/zenodo.23534954} and the Concept DOI
\href{https://doi.org/10.5281/zenodo.19210728}{10.5281/zenodo.19210728} are unchanged by it.

\subsection*{v4.9.0-r79 results}
\label{subsec:v490-r79-results}

\paragraph{Theoremlet v4.9.0-r79.1 (H67 confirmed: the four hidden conditions are gauge independent).}
Pre-registered.  Layer r72.1 showed that several algebraic gauges \(s=\Psi^\dagger M\Psi\) are transversal at the germ, and
that every transversal algebraic gauge gives the same symbol.  Whether the hidden conditions also agree depends on the
lower-order terms.
\begin{itemize}[leftmargin=2em]
\item \emph{Setup.}  For each gauge the high-precision fibre computation of r77.2 was repeated at \(q=6\) (70 digits, one
parallel runner per gauge).  The gauge series is \(w=2(\mathrm{Re}\,M\Psi,\mathrm{Im}\,M\Psi)\), with transversalities
\(|ds\cdot v|/|ds|\) at the germ in brackets:
  \begin{itemize}
  \item \(M=K_2\), the control (\(0.45\));
  \item \(M=1\), the density (\(0.077\));
  \item \(M=K_3\) (\(0.165\));
  \item \(M=Q_3\) (\(0.256\));
  \item \(M=I_3\otimes\beta\), the Lorentz scalar \(\bar\Psi\Psi\) (\(0.023\)).
  \end{itemize}
\item \emph{Result.}  In all five gauges the fibre of 6-jets is \(501-24=477\).  The Gram gaps are \(\ge4\times10^{56}\) and the
Newton residuals \(\le1.0\times10^{-51}\).  The polarization space has dimension 6 on 100 null directions, with gaps
\(\ge9\times10^5\), and dimension 0 off the cone.
\item \emph{Reading.}  The number of the four conditions and their support on the light cone do not depend on the gauge.  This
holds in particular for the covariant gauge \(\bar\Psi\Psi=S\).
\end{itemize}

\paragraph{Proof and measurement sources.}
r79.1 by \texttt{k167\_hp\_eichung.py}, run on GitHub runners with the workflow \texttt{hp-eichung.yml} and the data under
\texttt{ci/\allowbreak hp\_eich/}; the logs are \texttt{k167\_logs.txt}.  The note \texttt{Vorab\_\allowbreak Notiz\_\allowbreak H67}
and the result note lie in \texttt{notes/\allowbreak research/\allowbreak predictions/}.  \(\square\)

\paragraph{Consequence.}
Together with H66, the four hidden conditions belong to the covariant content of the comparison system, not to the gauge
fixing.  For all structural results so far (four conditions, formal integrability at the germ, six massless modes),
covariance is broken only by the prescribed background \((\Psi\text{ at the germ},A,h)\).  The next step is a native
dynamics for this background.  OP5 is not closed.

\subsection*{v4.9.0-r79 status in the three frames}
\label{subsec:v490-r79-frames}

\begin{itemize}[leftmargin=2em]
\item All statements concern the fixed OP5 comparison model in one chart, with one coframe reconstructed from \(\Psi\).  No
frame map M, T or P acts on them.
\end{itemize}
No statement of this layer concerns colour or charge.

\subsection*{v4.9.0-r79 register entries}
\label{subsec:v490-r79-register}

\begin{center}
\small
\begin{tabular}{@{}>{\raggedright\arraybackslash}p{0.34\textwidth}>{\raggedright\arraybackslash}p{0.58\textwidth}@{}}
\toprule
statement & status after v4.9.0-r79\\
\midrule
four hidden conditions & gauge independent in number and support (H67 confirmed, five gauges including \(\bar\Psi\Psi\))\\
source of Lorentz breaking & only the prescribed background \((\Psi,A,h)\)\\
next step & native dynamics of the background \(A,h\)\\
\bottomrule
\end{tabular}
\end{center}

\subsection*{v4.9.0-r79 ledger and guard update}
\label{subsec:v490-r79-ledger-guard}

\[
R_{\rm chk}^{v490\text{-}r79}=(1,0,0)
\]
for (gauge independence of the four conditions established; OP5 not closed).

The layer spends no guard:
\[
R_{\rm glob}^{(v490\text{-}r79,+)}=(0,1,1,1,1),
\qquad
\Gamma_{\rm guard}^{v490\text{-}r79}=(1,1,1,1,1),
\qquad
\beta_4=1/256\ \text{unspent}.
\]
It does not imply exact finite \(C5\) zero closure, exact all-readout terminality,
\(C6/T2\), a global photon-clock/RH theorem, or global all-readout closure.
